% Chapter 7 - Conclusiones

\chapter{Conclusiones}
\label{ch:conclusiones}

Este capítulo sintetiza los hallazgos experimentales, evalúa el cumplimiento de las hipótesis planteadas, identifica las contribuciones de la investigación, discute los resultados contraintuitivos, reconoce las limitaciones del estudio, y propone direcciones para trabajo futuro.

\section{Síntesis de Resultados}
\label{sec:sintesis_resultados}

El filtrado geométrico de muestras de texto generadas por modelos de lenguaje de gran escala mejora significativamente la clasificación de texto en escenarios de pocos ejemplos por clase. La evidencia proviene de más de 8,000 configuraciones experimentales evaluadas sobre 13 conjuntos de datos que abarcan 9 dominios textuales, 2 tareas de procesamiento de lenguaje natural y 5 modelos generativos de 4 proveedores distintos.

El resultado más fuerte se obtiene en la comparación con 10 métodos de aumentación sobre 1,890 configuraciones: la ponderación suave alcanza +2.61 puntos porcentuales respecto a SMOTE ($p < 0.0001$ con corrección de Bonferroni, $d = 0.95$, tasa de victoria del 88.9\%). El tamaño de efecto se clasifica como grande. En el experimento principal con 3,675 configuraciones, la mejora es de +2.25 puntos porcentuales ($p < 0.0001$, $d = 0.74$, tasa de victoria del 83.8\%).

El método funciona mejor con 10 a 25 ejemplos por clase, donde las ganancias alcanzan +4.52 y +1.88 puntos porcentuales respectivamente. Con 50 ejemplos, la mejora no alcanza significancia estadística. El rango óptimo de clases es de 2 a 20; con 77 o más clases el filtrado geométrico no aporta mejora.

La extensión a reconocimiento de entidades nombradas produce ganancias de +9.61 puntos porcentuales con tasa de victoria del 100\%, confirmando que los filtros geométricos son agnósticos a la tarea.

La evaluación con 5 modelos generativos (Gemini 3 Flash, GPT-5-mini, Claude 4.5 Haiku, Kimi K2.5, GLM-5) confirma que el método es agnóstico al LLM. Los 5 modelos producen mejoras significativas tras corrección de Bonferroni, con una dispersión de solo 0.82 puntos porcentuales entre el mejor (+2.31) y el peor (+1.49). El test de Friedman no detecta un efecto significativo del modelo generativo ($p = 0.086$), y la correlación inter-LLM ($\tau_K = 0.59$) indica acuerdo sustancial sobre qué configuraciones se benefician más del filtrado.

El análisis por clase individual revela que las clases difíciles (F1 base $<$ 30\%) se benefician desproporcionadamente, con ganancias medias de +10.44 puntos porcentuales, mientras que las clases fáciles (F1 $>$ 80\%) obtienen apenas +0.55 puntos porcentuales. Este patrón indica que el filtrado geométrico aporta mayor valor precisamente donde el clasificador tiene mayor margen de mejora: en las fronteras de decisión más inciertas del espacio de embeddings. En términos de F1 absoluto, la ponderación suave alcanza un macro F1 del 73.49\% frente al 71.24\% de SMOTE, contextualizando la mejora de +2.25 puntos porcentuales como una reducción del 7.8\% del error residual.

El experimento de aprendizaje curricular demuestra que la introducción ordenada de muestras sintéticas por score geométrico mejora el rendimiento respecto a la incorporación simultánea. El pico se alcanza al incluir el 50\% superior de los candidatos (+2.74 puntos porcentuales vs SMOTE, tasa de victoria del 86.2\%), superando al método estándar por +0.66 puntos porcentuales. Aproximadamente la mitad de las muestras generadas introduce ruido que degrada la calidad del clasificador.

La comparación entre filtrado geométrico y filtrado basado en modelo teacher confirma la superioridad del enfoque geométrico. El filtrado geométrico binario alcanza +2.78 puntos porcentuales frente a +2.43 del teacher binario. La diferencia es estadísticamente significativa ($p = 0.0001$), lo que demuestra que la distancia euclidiana al centroide de clase captura información más general que la confianza de un clasificador específico.

Las estadísticas geométricas del espacio de embeddings validan cuantitativamente el mecanismo del filtrado. Las muestras filtradas producen clústeres más compactos (distancia intra-clase: 0.822 vs 0.850 sin filtrar, $p < 0.001$) y mejor separados (distancia inter-clase: 0.616 vs 0.571, $p < 0.001$). El silhouette score mejora de 0.075 a 0.085 ($p = 0.029$), con victoria del filtrado en el 71\% de las configuraciones. Estos resultados confirman que el filtrado no solo selecciona muestras de mayor calidad sino que mejora la estructura geométrica del espacio de representación aumentado.

El experimento de aumentación híbrida adaptativa por clase demuestra que la especialización por dificultad no supera al filtrado geométrico uniforme. El filtrado puro alcanza +2.33 puntos porcentuales ($d = 0.93$, tasa de victoria del 91.0\%), mientras que la mejor estrategia híbrida obtiene solo +0.77 puntos porcentuales. El rendimiento mejora monótonamente a medida que más clases reciben filtrado geométrico en lugar de SMOTE, confirmando que la aplicación uniforme del mejor método es preferible a la especialización.

\section{Validación de Hipótesis}
\label{sec:validacion_hipotesis}

La \textbf{hipótesis principal H1} queda confirmada. El filtrado geométrico produce mejoras estadísticamente significativas respecto a SMOTE y todos los demás métodos de aumentación evaluados. La evidencia satisface los 3 criterios de aceptación: magnitud positiva (+2.25 puntos porcentuales), significancia estadística ($p < 0.0001$ con corrección de Bonferroni), y consistencia (tasa de victoria del 83.8\%).

La \textbf{hipótesis H1a} (filtrado) queda confirmada. El filtrado binario supera a la ausencia de aumentación por +2.08 puntos porcentuales ($p < 0.0001$, $d = 0.65$, tasa de victoria del 80.0\%). La selección geométrica de muestras aporta valor significativo respecto a la incorporación indiscriminada de todas las muestras generadas.

La \textbf{hipótesis H1b} (ponderación suave) queda parcialmente confirmada. En el experimento principal, la diferencia entre ponderación suave y filtrado binario es de +0.13 puntos porcentuales ($p = 0.10$), resultado que no alcanza significancia estadística. El estudio dedicado con hiperparámetros optimizados (normalización min-max, temperatura 0.5, selección top-N) muestra una ganancia de +3.66 puntos porcentuales con tasa de victoria del 84.5\%, pero este resultado representa un techo de rendimiento con configuración óptima, no el caso típico. La contribución principal del sistema propuesto reside en el filtrado geométrico en sí --- la selección de qué muestras incluir --- que produce ganancias estadísticamente significativas independientemente del esquema de ponderación. La ponderación suave aporta una mejora incremental cuya magnitud depende del ajuste de hiperparámetros, lo que limita su utilidad práctica como contribución autónoma.

La \textbf{hipótesis H1c} (generalización entre tareas) queda confirmada. Los filtros geométricos diseñados para clasificación de texto producen +9.61 puntos porcentuales en reconocimiento de entidades nombradas, con tasa de victoria del 100\% sobre 9 configuraciones. La adaptación no requiere modificación de los filtros, solo de la definición de clase.

La \textbf{hipótesis H1d} (efecto del clasificador) queda confirmada. El SVC lineal obtiene +2.98 puntos porcentuales ($d = 1.02$), Ridge +2.74 puntos porcentuales ($d = 1.04$), mientras que el bosque aleatorio obtiene +1.42 puntos porcentuales sin significancia estadística ($p = 0.11$). Los clasificadores lineales aprovechan más efectivamente las muestras filtradas.

\section{Contribuciones}
\label{sec:contribuciones}

Esta investigación aporta 9 contribuciones al campo de la aumentación de datos para clasificación de texto.

La primera contribución es un sistema de filtrado geométrico que selecciona muestras sintéticas de alta calidad en el espacio de embeddings. El sistema opera de manera desacoplada del modelo generativo, lo que se valida empíricamente con 5 LLMs de 4 proveedores distintos (Sección~\ref{sec:robustez_llm}). Un hallazgo relevante es que el filtro más simple (distancia euclidiana, cascada nivel 1) supera a los filtros multi-criterio complejos. El filtro combinado, que exige satisfacer múltiples criterios simultáneamente, produce resultados catastróficos con tasas de aceptación del 0\% en múltiples configuraciones.

La segunda contribución es el mecanismo de ponderación suave, que transforma las puntuaciones del filtro geométrico en pesos de muestra para el clasificador. La ganancia incremental sobre el filtrado binario es modesta en el experimento principal (+0.13 puntos porcentuales, $p = 0.10$), aunque alcanza +3.66 puntos porcentuales bajo configuración optimizada de hiperparámetros. Este resultado establece que la ponderación suave representa una extensión opcional cuyo beneficio depende del ajuste de temperatura y normalización, no un componente esencial del sistema. La contribución conceptual --- utilizar los scores geométricos como señal continua de calidad en lugar de un umbral binario --- permanece válida como dirección de investigación, aun cuando su impacto práctico sea incremental.

La tercera contribución es la validación sobre 10 conjuntos de datos que abarcan 7 dominios textuales distintos, con múltiples semillas aleatorias y corrección estadística rigurosa. La evaluación incluye pruebas t pareadas con corrección de Bonferroni, tamaños de efecto (d de Cohen), intervalos de confianza bootstrap al 95\% y tasas de victoria. El nivel de rigor estadístico supera al estándar habitual en la literatura de aumentación de datos.

La cuarta contribución es la demostración de generalización a reconocimiento de entidades nombradas. Los mismos filtros, sin modificación, producen mejoras de +9.61 puntos porcentuales en NER. Este resultado confirma que el filtrado basado en distancia euclidiana captura propiedades geométricas agnósticas a la tarea.

La quinta contribución es la evidencia comprensiva de que los filtros simples superan a los complejos. Esta observación, consistente a través de clasificación de texto, NER y 4 modelos de embeddings diferentes, constituye un hallazgo práctico relevante para la comunidad: la simplicidad en el diseño del filtro no solo reduce la complejidad de implementación sino que mejora el rendimiento.

La sexta contribución es la demostración de que el beneficio del filtrado geométrico se concentra en las clases más difíciles. El análisis desagregado por clase individual muestra una relación inversa entre el rendimiento base de la clase y la ganancia del filtrado: las clases con F1 inferior al 30\% obtienen ganancias medias de +10.44 puntos porcentuales, mientras que las clases con F1 superior al 80\% obtienen +0.55 puntos porcentuales. Este hallazgo tiene implicaciones prácticas directas: el filtrado geométrico es especialmente valioso para sistemas donde el rendimiento en clases difíciles o minoritarias es crítico.

La séptima contribución es la demostración mediante aprendizaje curricular de que la cantidad óptima de muestras sintéticas es inferior al total de candidatos disponibles. El rendimiento alcanza su pico al incluir el 50\% superior de los candidatos ordenados por score geométrico, con una ganancia de +0.66 puntos porcentuales sobre la incorporación simultánea de todos los candidatos. Este resultado proporciona una guía práctica para calibrar el presupuesto de muestras sintéticas y confirma que la calidad de la selección prevalece sobre la cantidad.

La octava contribución es la evidencia experimental de que la aplicación uniforme del filtrado geométrico supera a las estrategias híbridas adaptativas por clase. A pesar de que las clases difíciles se benefician desproporcionadamente del filtrado, especializar el método de aumentación por dificultad estimada degrada el rendimiento global. Este resultado, evaluado sobre 1,323 configuraciones con 3 métodos de estimación de dificultad, demuestra que la robustez del filtrado geométrico a través de todas las clases es preferible a la optimización local por subgrupos.

La novena contribución es la validación empírica de la independencia del modelo generativo. La evaluación con 5 LLMs de 4 proveedores (Gemini 3 Flash, GPT-5-mini, Claude 4.5 Haiku, Kimi K2.5, GLM-5) demuestra que el filtrado geométrico produce mejoras significativas independientemente del modelo utilizado. El test de Friedman ($p = 0.086$) no detecta diferencias significativas entre modelos, y la correlación inter-LLM ($\tau_K = 0.59$) confirma acuerdo sustancial. Este resultado transforma la afirmación de desacoplamiento, de un argumento de diseño a una evidencia empírica.

\section{Hallazgos Contraintuitivos}
\label{sec:hallazgos_contraintuitivos}

Cuatro resultados de esta investigación contradicen expectativas previas de la literatura.

El primer hallazgo contraintuitivo es que los filtros simples superan a los complejos. La cascada nivel 1 (solo distancia euclidiana) supera consistentemente al filtro combinado multi-criterio. Este último produce tasas de aceptación del 0\% en múltiples configuraciones, tanto en clasificación de texto como en NER. La Sección~\ref{sec:analisis_complejidad} desarrolla tres mecanismos que explican este resultado: la correlación matemática entre métricas en espacios normalizados (la distancia euclidiana y la similitud coseno producen rankings idénticos para vectores unitarios), la reducción de diversidad por intersección restrictiva de criterios, y la inestabilidad de las métricas de densidad local en el régimen de pocos ejemplos. Este análisis encuentra sustento teórico en Robinson et al. \cite{robinson2025manifold} y Du et al. \cite{du2025disentangling}. La simplicidad resulta ser una virtud, no una limitación.

El segundo hallazgo es que todas las líneas base modernas pierden contra SMOTE. La paráfrasis con T5, la aumentación contextual con BERT y el mixup en embeddings obtienen resultados negativos respecto a SMOTE simple. A pesar de su mayor sofisticación computacional y su uso de modelos preentrenados de gran escala, estas técnicas no generan diversidad genuina en el espacio de representación. Este resultado cuestiona la narrativa predominante de que las técnicas basadas en modelos preentrenados superan automáticamente a los métodos clásicos.

El tercer hallazgo es que la mejora iterativa del prompt no mejora las tasas de aceptación. El sistema de retroalimentación que analiza los motivos de rechazo y modifica el prompt en consecuencia produce tasas de aceptación decrecientes a lo largo de las iteraciones ($-12.3$ puntos porcentuales entre la iteración 0 y la 4). El cuello de botella no reside en la calidad del prompt sino en propiedades inherentes de la generación del LLM para ciertas distribuciones.

El cuarto hallazgo es que la especialización por dificultad de clase no mejora sobre la aplicación uniforme del filtrado geométrico. A pesar de que las clases difíciles obtienen ganancias 19 veces mayores que las fáciles (+10.44 vs +0.55 puntos porcentuales), reemplazar el filtrado geométrico por SMOTE en las clases fáciles degrada el rendimiento global. La mejor estrategia híbrida (+0.77 puntos porcentuales) queda muy por debajo del filtrado puro (+2.33 puntos porcentuales). El rendimiento mejora monótonamente con la proporción de clases que reciben filtrado geométrico, lo que indica que incluso las ganancias modestas en clases fáciles contribuyen al rendimiento agregado.

\section{Limitaciones}
\label{sec:limitaciones}

Esta investigación presenta varias limitaciones que delimitan el alcance de sus conclusiones.

La evaluación se realiza exclusivamente en inglés. El rendimiento del filtrado geométrico en otros idiomas no fue evaluado. Idiomas con menor representación en los modelos de embeddings podrían producir espacios de representación con propiedades geométricas diferentes.

La investigación utiliza un modelo principal de embeddings (all-mpnet-base-v2 de 768 dimensiones). Aunque la ablación con 4 modelos demuestra robustez (90.5\% de acuerdo mayoritario), todos los modelos pertenecen a la familia de transformers para inglés.

Existe una circularidad entre el espacio de filtrado y el de clasificación: el mismo modelo de embeddings se utiliza para calcular las distancias geométricas que determinan qué muestras se seleccionan y como representación de entrada para los clasificadores. Esto implica que el filtro optimiza la posición de las muestras en el mismo espacio donde el clasificador las evalúa, lo que podría favorecer artificialmente al filtrado geométrico. Tres líneas de evidencia mitigan esta preocupación. Primera, la ablación con 4 modelos de embeddings demuestra que cuando el filtrado se realiza con un modelo diferente al utilizado para clasificación, el acuerdo entre modelos es del 90.5\% (19 de 21 configuraciones), indicando que la selección geométrica captura propiedades de calidad que no son específicas del espacio de representación. Segunda, el filtrado geométrico iguala o supera al filtrado basado en modelo teacher, que utiliza un clasificador separado para evaluar la confianza de cada muestra. Tercera, los filtros generalizan a NER, donde la evaluación final emplea SpaCy y métricas seqeval a nivel de entidad, un pipeline completamente independiente del modelo de embeddings. No obstante, un diseño experimental ideal emplearía modelos de embeddings diferentes para filtrado y clasificación de manera sistemática. El 9.5\% de configuraciones sin acuerdo entre modelos (correspondientes a escenarios de 50 ejemplos, donde la aumentación aporta menos valor) sugiere que la circularidad podría tener efectos marginales en ciertos regímenes.

La evaluación multi-LLM emplea 5 modelos, todos de la generación más reciente de modelos comerciales. No se evalúan modelos de código abierto ejecutados localmente (como Llama o Mistral), donde la calidad de generación podría diferir. Adicionalmente, los filtros de contenido de los modelos evaluados impiden la generación para clases sensibles (discurso de odio, contenido religioso), lo que limita la aplicabilidad del método en dominios que involucren contenido potencialmente ofensivo. La clase ``hate\_speech'' fue rechazada por los 5 modelos en al menos una configuración.

Las generaciones del LLM se almacenan en caché, lo que produce varianza cero entre semillas para clasificadores determinísticos. El 80\% de las combinaciones método/clasificador presentan esta característica. Los intervalos de confianza podrían subestimar la varianza real que se observaría con regeneración completa de muestras en cada semilla. Restringiendo el análisis a clasificadores estocásticos, el resultado sigue siendo significativo (+1.79 puntos porcentuales, $p = 0.0015$, $d = 0.52$). Este análisis de sensibilidad constituye la prueba más conservadora de las hipótesis de esta investigación. Un diseño experimental alternativo regeneraría las muestras sintéticas con diferentes prompts o temperaturas del LLM en cada semilla, proporcionando una estimación más fiel de la variabilidad real. La significancia del resultado con clasificadores estocásticos ($d = 0.52$, efecto mediano) sugiere que las conclusiones son robustas, pero los intervalos de confianza del análisis completo deben interpretarse como cotas inferiores de la incertidumbre real.

El resultado a 50 ejemplos por clase no alcanza significancia estadística ($p = 0.16$). El método es más valioso en escenarios de recursos limitados (10-25 ejemplos). Con datos de entrenamiento moderados, la aumentación aporta valor marginal decreciente.

La evaluación de NER emplea SpaCy como modelo de reconocimiento de entidades. Los modelos de NER de última generación podrían producir resultados diferentes. La elección de SpaCy responde a su reproducibilidad y disponibilidad, pero limita la comparabilidad con la literatura más reciente de NER.

El número de clases limita las ganancias. Con 77 clases (banking77) el resultado es negativo ($-0.99$ puntos porcentuales) y con 150 clases (clinc150) es neutro (+0.01 puntos porcentuales). El filtrado geométrico basado en distancia pierde efectividad cuando el espacio de embeddings se fragmenta excesivamente.

\section{Trabajo Futuro}
\label{sec:trabajo_futuro}

La extensión multilingüe representa la dirección más inmediata. Evaluar el filtrado geométrico sobre modelos de embeddings multilingües permitiría verificar si las propiedades geométricas que fundamentan el filtrado se mantienen en otros idiomas.

La evaluación con modelos generativos de código abierto ejecutados localmente constituye una segunda línea prioritaria. La validación multi-LLM de esta investigación emplea modelos comerciales con API; evaluar modelos como Llama y Mistral ejecutados sin filtros de contenido permitiría verificar si el método funciona en dominios sensibles donde los modelos comerciales rechazan la generación.

La búsqueda conjunta de hiperparámetros mediante optimización bayesiana podría revelar configuraciones superiores. La optimización secuencial utilizada en esta investigación explora cada parámetro independientemente. Una búsqueda conjunta capturaría interacciones entre el modelo de embeddings, el nivel de filtrado, la temperatura de ponderación y el número de candidatos generados.

La selección adaptativa de temperatura por clase ofrece potencial de mejora. Clases con alta variabilidad interna podrían beneficiarse de temperaturas más altas en la generación, mientras que clases compactas podrían preferir generaciones más conservadoras.

El presupuesto adaptativo de generación por clase asignaría más recursos a clases difíciles. En lugar de generar el mismo número de candidatos para cada clase, un sistema adaptativo invertiría más en las clases con menor tasa de aceptación o mayor dificultad de clasificación.

La estimación formal de ratios de densidad permitiría formalizar el muestreo por importancia que subyace a la ponderación suave. Una estimación explícita de la relación entre la distribución de generación del LLM y la distribución real de cada clase proporcionaría pesos más fundamentados teóricamente.

La extensión a otras tareas de procesamiento de lenguaje natural, como generación de preguntas, resumen extractivo y traducción automática, verificaría el alcance de la generalización demostrada en esta investigación entre clasificación y NER.
